\chapter{Density Functional Theory}
\label{sec:DFT}

The density functional theory (DFT) is recognized as a overwhelmingly successful method in quantum chemistry and the overall study and understanding of matter. As illustrated in figure 5.1 this is a tremendously popular method with rapid growth to this day due to improvements to both the method and computational power. \textbf{More elaborate introduction, write later.}
\begin{figure}[H]
\includegraphics[width=\textwidth]{theory/dftScope.png}
\caption{Number of DFT studies per year from 1980 to 2021. \textbf{Add reference?}}
\end{figure}

The concept of the density functional theory is to efficiently solve the problems posed by quantum physics, most notably the many-particle Shr\"{o}dinger equation. Thus in this section we will begin by reviewing some central concepts of quantum physics before presenting the main Hohenberg-Kohn formulation of DFT. The content in this part is inspired from the lecture notes [REF] written by Clas Persson in the course FYS-MENA4111 at the University of Oslo.
\textbf{References for theory: fys-mena4111 lecture notes and book}

Quantum mechanics is the engine of density functional theory and to truly understand the method and it's limitations a knowledge of the engines is necessary. 
\section{Review of Quantum Mechanics}

\subsection{The Shr\"{o}dinger equation}
All theoretical information of a material can be derived from the Shrodinger equation, which came to the world in 1926. Since this the major question how been how this equation can be solved. 
The time-dependent Schr\"{o}dinger equation for one electron is given bellow

\begin{equation}
    i\hbar\frac{\partial}{\partial t}\Psi(\vec{r}, t) = \hat{H}(\vec{r},t)\Psi(\vec{r}, t).
\end{equation}

This equation consists of the the wavefunction $\Psi(\vec{r},t)$ which describe the electron, and the Hamiltonian $\hat{H}(\vec{r},t)$ where $\vec{r}$ and $t$ is the spatial position and time respectively. The Hamiltonian describes the systems total energy by a kinetic part $T = \frac{-\hbar^2\nabla^2}{2m_e}$, where $m_e$ is the electron mass and $\hbar$ is the constant, and a potential energy operator U, typically an external potential denoted $V_{ext}(\vec{r}, t)$. Wavefunctions of stationary states are eigenfunctions of the Hamiltionian and are denoted as $\psi_{\kappa}(\vec{r}, t)$ with an energy eigenvalue $\epsilon_{\kappa}$. The subscript $\kappa$ indicate that the wavefunction correspond to the $\kappa$ eigenstate. Above we included the time-dependent Shrodinger equation, but almost all cases involving Quantom physics employ rather the time-independent Shroidinger eqaution in which the external potential is independent of time. Bellow we show the time-independent Shrodinger equation for the eigenvalues $E_k$ of the $k$-th eigenfunction $\psi_k(\vec{r})$

\textbf{Talk about general wavefunctions vs eigenfunctions or skip this?)}

\begin{equation}
	\left(-\frac{\hbar^2\nabla^2}{2m_e} + V_{ext}(\vec{r}) \right)\psi_{\kappa}(\vec{r}) = E_k \psi_k(\vec{r})
\end{equation}

Solving the single electron time-independent equation often results in infinite eqienstates that the electron can occupy. The most probable state the electron occupy is the lowest energy state called the ground state, this state is indicated by $\kappa = 0$. Extending to a system comprised of multiple particles we have the many-body Shr\"{o}dinger equation with the many-body wavefunction 

The many-body wavefunction is defined as
\begin{equation}
\Psi^{en}(\vec{r}, \vec{R}) = \Psi^{en}(\vec{r_1}, \vec{r_2}, \dots \vec{r_{N_e}}, \vec{R_1}, \vec{R_2}, \dots \vec{R_{N_n}}),
\end{equation}
 where $r_j$ denote the coordinates of the j:th electron and likewise $R_{\alpha}$ describe the coordinates of the $\alpha$:th nucleus, and the subscript "en" subscript means that we consider both the electrons and nuclei. Furthermore we have the many-body Hamiltonian $H^{en}$, composed of the kinetic energy of $N_e$ electrons $T_e$, the interaction energy between electrons $U_{ee}$, the kinetic energy of $N_n$ nuclei, the coulomb interaction between nuclei $U_{nn}$, and finally the attractive interaction between nuclei and electrons $U_{en}$. In the many-body Hamiltonian the following notations are used $m_n$ = nuclei mass, $q$ = particle charge, $\alpha$ = nuclei number, $Z_a$ = atom number of nuclei $\alpha$. In final we get 
\begin{align}
    \hat{H} &= T_e + T_n + U_{ee} + U_{nn} + U_{en} \\
    \begin{split}
        &= -\sum_{j=1}^{N_e}\frac{\hbar^2\nabla_j ^2}{2m_e} - \sum_{\alpha=1}^{N_n}\frac{\hbar^2 \nabla_\alpha ^2}{2m_n} + \sum_{j=1}^{N_e}\sum_{j'<j}\frac{q^2}{|r_j - r_{j'}|} \\
    &+ \sum_{\alpha=1}^{N_n}\sum_{\alpha' < \alpha}\frac{q^2Z_\alpha Z_{\alpha'}}{R_\alpha - R_{\alpha'}} -\sum_{j=1}^{N_e}\sum_{\alpha=1}^{N_n}\frac{q^2Z_\alpha}{|r_j - R_\alpha|},
    \end{split}
\end{align} 
 
and the many body Schrodinger equation with total energy eigenvalue $E_{\kappa}^{en}$ of the whole system in eigenstate $\kappa$
 
\begin{equation}
H^{en}\Psi^{en}_{\kappa}(\vec{r}, \vec{R}) = E_{\kappa}^{en}\Psi_{\kappa}^{en}(\vec{r}, \vec{R}).
\end{equation}

The first step to solving the many-body problem is to obtain an exact expression for the many particle wavefunction $\Psi^{en}(\vec{r}, \vec{R})$ in order to apply the single-particle operators in equation (5.4). If we consider a simplified system consisting of just two electrons, the problem is reduced to finding $\Psi_{\kappa}(\vec{r_1}, r_2)$ that is a function of $psi_1(\vec{r_1})$ and $\psi_2(\vec{r_2})$. In the Hartree approach this is solved by considering the two electrons to be independent of each-other, and employ variable separation to express the two particle wavefunction as

\begin{equation}
\Psi_{\kappa}(\vec{r_1}, r_2) = psi_1(\vec{r_1})psi_2(\vec{r_2}).
\end{equation}

The limitation of the Hartree approach is that the electrons, which are fermions, in this formulation are distinguishable and hence does not obey the Pauli exclusion principle of fermions. This is corrected by the Harte-Fock approximation which introduces a spin function $\chi_{mp}(s_1, s_2)$ to equation (5.4) to make it anti-symmetric with respect to the particle coordinates. The Hartree-Fock approximation is expressed as

\begin{equation}
\Psi_{\kappa}(\vec{r}_1, \vec{r}_2) = \frac{1}{2} \left\{ \psi_1(\vec{r}_2) \psi_2(\vec{r}_2) \pm \psi_1(\vec{r}_2) psi_2(\vec{r}_1) \right\} \chi_{\mp}(s_1, s_2).
\end{equation}  

The difference in energy from the improved wavefunction in Harte-Fock compared to the Hartee approximation is called the exchange energy. Note however that Hartee-Fock is not a complete description either as it fails to model the electron correlations.  

\subsection{Simplifications and approximations to solve the many-electron Shr\"{o}dinger equation}

\subsubsection{Born-Oppenheimer}
Challenges with solving many-particle Shr\"{o}dinger equation is i) computationally expensive, ii) need to know how $\Psi$ depends on single particle wavefunctions $\psi_k$. To solve this complex problem, we need approximations. Particularly Born-Oppenheimer and Harte-Fock approximations. The first makes the cleaver and reasonable assumption that since the electron mass is negligibly small in comparison to that of a nuclei, we can treat the nuclei as point charges, enabling us to divide the eigenfunction into a separate electronic and nuclear part, ie
\begin{equation}
    \Psi_{k}^{en}(\vec{r}, \vec{R}) \approx \Psi_k(\vec{r}, \vec{R}) \Theta_k(\vec{R})
\end{equation}
where we have written the complete wavefunction in terms of an electronic part $\Psi_k(\vec{r}, \vec{R})$ and nuclear part $\Theta_k(\vec{R})$. The dependencies come from the fact that electrons can respond instantaneously to new positions of the nuclei, therefore the $\vec{R}$ dependence. Writing this in terms of the Hamiltonian we get
\begin{align}
    &\left( T_{e} + U_{ee} + U_{en} \right) \Psi_k(\vec{r}, \vec{R}) = E_k(\vec{R})\Psi_k(\vec{r}, \vec{R}) \\
    &\left( T_{n} + U_{nn} + E_k(\vec{R}) \right) \Theta(\vec{R}) = E_{k}^{en}(\vec{R})\Theta_k(\vec{r}, \vec{R}).
\end{align}
The two sections are interrelated through the electronic energy eigenvalue $E_k(\vec{R})$. Furthermore, the left hand side of the nuclear part can be simplified to $U_{nn} + E_k(\vec{R})$, assuming that the kinetic energy of point charges is zero. This simplified expression for the nuclear left hand side is called for the potential energy surface (EPS).

\subsubsection{Hartree-Fock}
The next step in line is to find a wavefunction that can describe all electrons in a system. This was originally done by Hartree, which assumed that electrons can be described independently and suggested the ansatz for a two-electron wavefunction
\begin{equation}
    \Psi_k(\vec{r_1}, \vec{r_2}) = A \cdot \psi_1(\vec{r_1}) \psi_2(\vec{r_2}),
\end{equation}
where A is the normalization constant. However this approximation does not account for the fact that electrons are indistinguishable and hence does not obey the Pauli exclusion principle for ferminons. This was overcome with the Hartree-Fock approximation that implement an anti-symmetric wavefunction. The full expression is given bellow
\begin{equation}
    \Psi_k(\vec{r_1}, \vec{r_2}) = \frac{1}{\sqrt{2}} \Big( \psi_1(\vec{r_1}) \psi_2(\vec{r_2}) - \psi_1(\vec{r_2})\psi_2(\vec{r_1}) \Big)
\end{equation}

The Hartree-Fock (HF) approximation makes the electrons distinguishable and hence obey the Pauli exclusion principle, this means that the exchange energy is accounted for. On the other side, HF is not a complete description as it fails to model the electron correlations. 

\subsubsection{The Variational principle}
In materials science, the overarching concern is the ground-state properties of a system. This can be found efficiently and easy by whats known as the variational principle. This states that the energy of any trial wavefunction will always be higher than the ground-state energy $E_0$, ie
\begin{equation}
    E_0 = \langle\psi_0|H|\psi_0\rangle \leq \langle\psi|H|\psi\rangle = E
\end{equation}
This enable us to find the ground state energy and corresponding wavefunction by a minimization technique. Next, we will present the basics of the density functional theory for how these equations can be solved numerically and efficiently in order to study real materials and systems.

\section{Fundamentals of Density-Functional Theory}
The density functional theory was developed by Hohenberg and Kohn in 1964 and revolved around the fact that the ground-state density can be expressed in terms of the ground-state wavefunction. We have
\begin{equation}
    n_0(r) = |\Psi_0(r)|,
\end{equation}
furthermore the theorem states that all ground-state physical properties can be found as unique functionals of the ground-state density. The biggest upside of this, is that instead of trying to solve the many-body Shr\"{o}dinger equation to obtain the ground-state wavefunction, we have reduced the computational complexity from $3N_e$ to $3$. Thus, the Hohenberg and Kohn density functional theory makes for a promising and effective method to obtain the ground-state properties of a system, given that the exact electron density functional is known. However, this is still 60 years later unknown. 

The density functional theory build on two specific theories, called the Hohenberg-Kohn theorems. They are:
\begin{enumerate}
    \item For any system of interacting particles in an external potential $V_{ext}$, the density is uniquely determined.
    \item There exists a variational principle for the energy density functional such that, if $n$ is not the electron density of the ground-state, them $E[n_0] < E[n]$.
\end{enumerate}
The proof behind both theorems can be found in appendix .. A direct result of the second theorem is the energy can be described as a function of the density
\begin{equation}
    E[n] = T[n] + U_{ee}[n] + U_{en}[n],
\end{equation}
where the first two terms $T[n]$ and $U_{ee}[n]$ make up the Hohenberg-Kohn functional. 

We know move on to the Kohn-Sham equations, in which Kohn and Sham expressed the exact ground-state density from Hartree type wavefunctions. 
\begin{equation}
    \Psi(\vec{r}_1, \vec{r}_2 , .., \vec{r}{N_e}) = \psi_1^{KS}(\vec{r}_1)\psi_2^{KS}(\vec{r}_2)...\psi_{N_e}^{KS}(\vec{r}_{N_e})
\end{equation}
In which, $\psi_j^{KS}$ are auxiliary independent single-particle wavefunctions. We know modify the equation for total energy as a function of density defined by the second theorem, to include the single auxiliary wavefunctions and their corresponding kinetic energy and interaction energy. We get:
\begin{equation}
    E[n] = T_s[n] + U_s[n] + U_{en}[n] + (T[n] - T_s[n]) + (U_{ee}[n] - U_s[n]).
\end{equation}
with the s subscript denoting the single particle wavefunctions. The latter two terms are known as the exchange-correlation energy $E_{xc}$
\begin{equation}
    E_{xc}[n] = \Delta T + \Delta U
\end{equation}

This term is responsible for the many-electron interaction. The complete total energy functional can now be expressed as
\begin{equation}
    \begin{split}
    E[n] &= \overbrace{\sum_j \int \psi_j^{KS*}\frac{-\hbar^2 \nabla^2}{2m}\psi_j^{KS}d\vec{r}}^{T_s[n]} + \overbrace{\frac{1}{2}\frac{1}{4\pi \epsilon_0} \int \int q^2 \frac{n(\vec{r})n(\vec{r}')}{|\vec{r} - \vec{r}'|}d\vec{r}\vec{r}'}^{U_s[n]} \\ 
        &+ \underbrace{\int V_{en}(\vec{r}n(\vec{r})d\vec{r}}_{U_en[n]} + \underbrace{(T[n] - T_s[n]) + (U_{ee}[n] - U_s[n])}_{E_{xc}[n]}
    \end{split}
\end{equation}

Finally we write the complete expression for the Kohn-Sham single-electron equations given an exact exchange-correlation energy and utilizing the variational principle described previously
\begin{equation}
    \bigg\{ -\frac{\hbar^2}{2m_e}\nabla^2_s + v_H(\vec{r}) + V_{en}(\vec{r}) + V_{xc}(\vec{r} \bigg\}\psi_s^{KS}(\vec{r}) = \epsilon_s^{KS}(\vec{r})\psi_s^{KS}(\vec{r}),
\end{equation}
Define $V_H$, and $V_{xc}$ and mention that the former include self interaction that can be accounted for in XC functional. Finally, the total energy of the many-electron system is defined as
\begin{equation}
    E[n] = \sum_j \epsilon_j^{KS} - \frac{1}{2}\frac{1}{4\pi\epsilon_0} \int \int q^2 \frac{n(\vec{r})n(\vec{r}')}{|\vec{r} - \vec{r}'|}d\vec{r}d\vec{r}' + E_{xc}[n] - \int V_{xc}(\vec{r})n(\vec{r})d\vec{r}.
\end{equation}
This is the fundamental working principle of the density functional theory and Kohn-Sham equations.

\section{Limitations of DFT}
\begin{itemize}
    \item Local minima method
    \item Not exact $V_{xc}$, means we must compromise between accuracy and cost, and choose between the different methods for specific application. There is no one best overall method that is superior for all purposes. 
    \item kohn-sham eigenfunctions are not the exact eigenfunctions. 
    \item DFT in its original formulation is only valid for the ground state. Thus excited state application of DFT, all though very possible have a lesser theoretical footing.
    \item The band-gap calculation is also complicated by the self-interaction error which arises in the occupied states in standard DFT, and in the unoccupied states in Hartree-Fock. If you consider the "true" band-structure, then semi-local DFT has a spurious self-interaction in the occupied states, which over-delocalises them and forces them up in energy, thus reducing the band-gap; Hartree-Fock has the same problem but for the *unoccupied* states, so they are over-delocalised and forced up in energy, which increases the band-gap. \cite{bandgap_dft_forum}
    \item References on the band gap with DFT based methods \cite{bandgap_dft_1}, \cite{bandgap_dft_2}
\end{itemize}

